\section{Kritische Reflexion und Grenzen}

Ein zentrales Problem war, dass die verfügbaren Ressourcen nicht zu Beginn des Projekts geklärt wurden. Während der Implementierung sind Einschränkungen aufgetreten, die bei sorgfältiger Vorplanung erkennbar gewesen wären. Eine Architektur auf Dienste oder Zugänge auszulegen, die später nicht oder nur eingeschränkt nutzbar sind, führt zu vermeidbarem Mehraufwand.

Ein solches Problem tritt auch in der Praxis häufig auf, da bei der Planung oft
stärker berücksichtigt wird, was technisch grundsätzlich möglich wäre, als was
mit den tatsächlich verfügbaren Mitteln realistisch umgesetzt werden kann. Gerade
bei Schulprojekten, in denen Zugänge zu Cloud-Diensten oder Lizenzen nicht
selbstverständlich vorhanden sind, hätte diese Klärung früher erfolgen müssen.

Positiv war die Zusammenarbeit im Team. Auftretende Probleme wurden gemeinsam bearbeitet und gelöst, nicht ignoriert oder aufgeschoben. Die Zusammenarbeit funktionierte auch bei unvollständiger Vorplanung. Probleme wurden dadurch nicht verhindert, aber zügig behoben.

Die Empfehlung für künftige Projekte: Ressourcenplanung und technische Planung sollten gemeinsam am Anfang stehen. Die Architektur sollte erst konkretisiert werden, wenn die verfügbaren Mittel bekannt sind.
